\section{Pila del sistema de generación: DSL, Context, Agent/Workflow y System1}

Después de discutir por qué hacer video, esta sección entra en la metodología central que propone Wang Guan (王冠): el sistema de generación. Un sistema de generación no es un solo modelo ni un solo workflow, sino un método de producción organizado en torno a cierto tipo de bien de información. Integra en un sistema la expresión del problema, el context efectivo, el agent/workflow, las capacidades del modelo y la evaluación.

\begin{figure}[H]
\centering
\includegraphics[width=0.92\textwidth]{generation-system-stack.png}
\caption{Pila del sistema de generación: System1, DSL, Context, Agent/Workflow y evaluación trabajan en conjunto. Diagrama conceptual propio, elaborado a partir de la conversación de 02:01:44--02:25:49.}
\end{figure}

\begin{knowledgebox}{Lectura de la figura: el sistema de generación es principalmente System2}
System1 es el modelo base; el DSL sirve para expresar el problema y el método de producción; el Context aporta tokens efectivos; el Agent/Workflow organiza los pasos; la evaluación mide el grado de inteligencia del sistema. Wang Guan considera que agent y workflow no tienen por qué oponerse: en esencia, ambos producen un context más efectivo para que System1 genere mejores salidas.
\end{knowledgebox}

\subsection{DSL: expresar el método de producción}

DSL es Domain-Specific Language, lenguaje específico de dominio. Aquí no se refiere necesariamente a un lenguaje de programación tradicional, sino al sistema de expresión de una tarea de generación. En el caso del video, el DSL debe poder describir tomas, ritmo, imagen, estilo, material, transiciones, público objetivo y método de producción. Sin DSL, el sistema solo puede acumular prompts; con DSL, el sistema puede reutilizar el método de producción de forma estable.

\subsection{Context Engineering: el nuevo campo principal del gerente de producto}

Esta sección pasa del DSL al context, porque el resultado de un sistema de generación no depende solo del modelo, sino también de lo que el modelo ve antes de generar y de cómo entiende la estructura de la tarea.

Wang Guan considera muchos agents, workflows, bases de conocimiento y repositorios especializados como mecanismos para generar context. La salida del modelo depende de los tokens previos; lo que el producto debe hacer es convertir el know-how del negocio en un context que el modelo pueda entender, que sea económicamente rentable y que dé resultados estables. Este sistema no es necesariamente el más complejo en lo técnico, pero depende en gran medida de la comprensión del negocio y del criterio de producto.

\begin{warningbox}{No subestimar el Workflow como algo «sin contenido técnico»}
Workflow, Agent, repositorios especializados y DSL sirven para proporcionar al modelo un context efectivo. Su valor no reside solo en la complejidad de ingeniería, sino en si logran expresar de forma estable los métodos del negocio, la estructura de la tarea y los objetivos de evaluación.
\end{warningbox}

\subsection{Métrica de estrella polar: el grado de inteligencia del sistema}

Esta sección responde «qué optimiza en realidad un producto de IA». El valor comercial es el mínimo indispensable, pero si solo se mira el ingreso a corto plazo, se pierde de vista el crecimiento del sistema de generación como capacidad de producción en sí misma.

Wang Guan sostiene que un producto AI Native debe tener valor comercial, desde luego, pero que la métrica de estrella polar de nivel superior es el «grado de inteligencia del sistema». No se trata de algo esotérico, sino de preguntar: ¿el sistema puede producir de forma automática, elegir métodos de forma automática, corregirse de forma automática y llegar por sí mismo a mejores salidas? En un sistema de generación, el usuario no solo compra una herramienta, sino que usa un sistema de producción que cada vez sabe hacer mejor su trabajo.

\begin{figure}[H]
\centering
\includegraphics[width=0.96\textwidth]{ai-product-north-star.png}
\caption{Estrella polar de un producto de IA: por encima del valor comercial, medir el producto por el grado de inteligencia del sistema. Diagrama conceptual propio, elaborado a partir de la conversación de 02:50:34--03:05:45.}
\end{figure}

\begin{knowledgebox}{Lectura de la figura: el grado de inteligencia no es un lema abstracto}
La figura va del valor comercial al diseño de capacidades, y de ahí al context, la producción automática y el grado de inteligencia. El grado de inteligencia del sistema puede manifestarse en: si entiende el objetivo, si elige el método adecuado, si reduce los pasos manuales, si puede mejorar a partir de sus fallas y si permite al usuario lograr una producción de mayor calidad a menor costo.
\end{knowledgebox}

\subsection{Resumen de la sección}

El sistema de generación es el núcleo con el que una empresa de aplicaciones puede construir una ventaja defensiva. Convierte las capacidades del modelo en métodos de producción, el know-how del negocio en context y cada generación aislada en una recipe reutilizable, y mide el avance del producto por el grado de inteligencia del sistema.
